%
% zweipunkte.tex -- Erweiterung einer holomorphen Funktion mithilfe 
%                   einer Potenzreihe
%
% (c) 2021 Prof Dr Andreas Müller, OST Ostschweizer Fachhochschule
%
\documentclass[tikz]{standalone}
\usepackage{amsmath}
\usepackage{times}
\usepackage{txfonts}
\usepackage{pgfplots}
\usepackage{csvsimple}
\definecolor{darkred}{rgb}{0.8,0,0}
\definecolor{darkgreen}{rgb}{0,0.6,0}
\def\r{2.1}
\usetikzlibrary{arrows,intersections,math,calc}
\begin{document}
\def\skala{1}
\def\punkt#1#2{
	\fill[color=white] #1 circle[radius=0.05];
	\draw[color=#2] #1 circle[radius=0.05];
}
\begin{tikzpicture}[>=latex,thick,scale=\skala]

\coordinate (A) at (3,2.4);
\coordinate (B) at (4.5,3);

\fill[color=gray!20,opacity=0.5] (A) circle[radius=\r];
\fill[color=gray!20,opacity=0.5] (B) circle[radius=1.4];

\begin{scope}
	\clip (A) circle[radius=\r];
	\fill[color=darkred!20,opacity=0.5] (B) circle[radius=1.4];
\end{scope}

\begin{scope}
	\clip (B) circle[radius=1.4];
	\fill[color=darkgreen!20,opacity=0.5]
		($(A)+(-60:\r)$) arc(-60:90:\r)
		--(6.8,4.5) -- (6.8,0) -- cycle;
\end{scope}

\node at ($(A)+(-90:\r)$) [above] {$D_1$};
\node at ($(B)+(60:1.5)$) [below left] {$D_2$};

\draw[->,line width=0.5pt] (A) -- ++(120:2.1);
\draw[->,line width=0.5pt] (B) -- ++(-10:1.4);

\node at ($(A)+(120:1.4)$) [left] {$\varrho_1$};
\node at ($(B)+(-10:1.1)$) [above] {$\varrho_2$};

\punkt{(A)}{black}
\punkt{(B)}{black}

\node at (A) [below] {$z_1$};
\node at (B) [above] {$z_2$};


\draw[->] (-0.1,0) -- (6.8,0) coordinate[label={$\operatorname{Re}$}];
\draw[->] (0,-0.1) -- (0,5.3) coordinate[label={right:$\operatorname{Im}$}];

\end{tikzpicture}
\end{document}

